\subsection{Fasting and Almsgiving}\label{sec:fasting}

St Leo the Great opened one of his sermons on the December fast with a reason that belongs to the
season of the year:

\begin{quote}
Dearly beloved brethren, with the care which becometh us as the shepherd of your souls, we urge
upon you the rigid observance of this December Fast. The month of December hath come round again,
and with it this devout custom of the Church. The fruits of the year, which is drawing to a close,
are now all gathered in, and we most meetly offer our abstinence to God as a sacrifice of
thanksgiving.\par
{\smallskip\noindent\small ---Leo the Great, \textit{Sermo} 13, tr. Bute, \textit{The Roman
Breviary} (1908), Winter, Third Sunday of Advent, lesson iv, as read before 1960.\par}
\end{quote}

\noindent In another sermon Leo joins the fast to almsgiving and to prayer: ``For forgiveness
of sins is most efficaciously prayed for with almsgiving and fasting, and supplications that are
winged by such aids mount swiftly to God's ears'' (\textit{Sermo} XVI.2, tr. Feltoe). The Roman
Catechism gives the same teaching in its treatment of prayer:

\begin{quote}
To prayer let us unite fasting and alms-deeds. Fasting is most intimately connected with prayer:
When cloyed with meat and drink, the mind is so pressed down as not to be able to raise itself to
the contemplation of God, or comprehend the utility of prayer. Alms-deeds have also an intimate
connexion with prayer.\par
{\smallskip\noindent\small ---\textit{Catechism of the Council of Trent}, tr. J. Donovan
(Baltimore: Lucas Brothers), p.~332.\par}
\end{quote}

\noindent In its treatment of penance the same Catechism brings all satisfaction under these three
heads, prayer, fasting and almsgiving.

\subsubsection{What the law required, and what it requires}

The law on fasting has changed more than once since the books of 1962, and what bound then is
easily confused with what binds in 2026.

\textit{The Latin Church, 1918 to 1983.} The Code of Canon Law of 1917, in force from 19 May 1918
to 26 November 1983, required fast and abstinence on the Ember days and on the Vigil of the
Nativity (can.~1252 \S\,2). It did not name Advent as a time of fasting. When a day of fast fell on
a Sunday the law ceased for that day, and vigils were not moved earlier (can.~1252 \S\,4). A
decree of the Sacred Congregation of the Council of 3 December 1959 gave all the faithful leave to
keep the fast and abstinence of the Vigil on 23 December instead of the 24th.

\textit{The United States in 1889.} The \textit{Manual of Prayers} issued for the United States in
1889 lists the Fridays of Advent among the fasting days, with the Ember days and the Vigil of
Christmas, and notes that in some dioceses, ``as in the Province of New Orleans,'' the Fridays of
Advent are not fasting days. It also says that a fast falling on a Sunday is kept on the Saturday
before, and it presents these practices as ``enjoined by laws binding strictly on the conscience.''

\textit{England in 1870.} The English translator of Guéranger's \textit{Liturgical Year}, Dom
Laurence Shepherd, added a note to Guéranger's history of the season: ``Our English observance of
Fast and Abstinence on the Wednesdays and Fridays in Advent, may, in some sense, be regarded as a
remnant of the ancient discipline.'' His note attests a custom observed in England in his day.

\textit{The Latin Church in 2026.} Under the Code of 1983, in force from 27 November 1983 and current
on the First Sunday of Advent 2026, the penitential days and times of the universal Church are
every Friday of the year and the season of Lent (can.~1250). Abstinence is kept on every Friday
unless it is a solemnity, and fast and abstinence on Ash Wednesday and Good Friday (can.~1251).
Abstinence binds those who have completed their fourteenth year; fasting binds all who have come
of age until the beginning of their sixtieth year (can.~1252). A conference of bishops may determine these more
precisely and may substitute other forms of penance, especially works of charity and exercises of
piety (can.~1253). Paul VI's \textit{Paenitemini} of 17 February 1966 had already named Lent and
the Fridays of the year, with Ash Wednesday, as the days of penance (norm II), and had given
conferences the same power to substitute (norm VI). Neither the Code nor \textit{Paenitemini}
names Advent.

\textit{The United States in 2026.} The bishops of the United States, in their pastoral statement of
18 November 1966, effective on the First Sunday of Advent, 27 November 1966, asked Catholics ``to
make the Advent season, beginning with 1966, a time of meditation on the lessons taught by the
liturgy and of increased participation in the liturgical rites by which the Advent mysteries are
exemplified and their sanctifying effect is accomplished'' (par.~7). They imposed no Advent fast.
They stated that ``Vigils and Ember Days, as most now know, no longer oblige to fast and
abstinence,'' while leaving the devout free to prepare for each festival ``by a day of particular
self-denial, penitential prayer and fasting'' (par.~17). Friday remains a penitential day; the
statement ended abstinence on the Fridays outside Lent as an obligation under pain of sin, but
gave ``first place to abstinence from flesh meat'' among the ways of keeping Friday (pars.~22--25).
Under the Code of 1983 the bishops' complementary norm keeps the 1966 norms in force, and an
amendment of 1984 sets the age of fasting ``from the completion of the eighteenth year to the
beginning of the sixtieth''.

In short: for a Catholic of the Latin Church in the United States in Advent 2026, every Friday is
a penitential day, and neither the universal law nor the American norms make any day of Advent a
day of obligatory fast. The Ember days and the Vigil, which bound under the Code of 1917, bind
under neither.

\subsubsection{A fast of one's own}

The book's counsel is to keep some fast in Advent all the same, measured and modest. The bishops
of 1966, imposing no fast before any feast, left the devout free to prepare for each festival by
fasting (par.~17). A household might keep the Ember days with one full meal,
or the Fridays of Advent with a plainer table, as the English custom kept the Wednesdays and
Fridays. Health, age, work, pregnancy and the care of others all set limits, and a confessor or
director is the right judge of them. A fast that makes one irritable with one's family has missed
Leo's point, which was that the mind should fast from vices as the body fasts from food.

\subsubsection{Alms}

Leo joined alms to the December fast. In the sermon that names the four seasonal fasts he
writes:

\begin{quote}
And hence we warn you, beloved, in fatherly affection, to make this winter fast fruitful to
yourselves by bounteous alms, rejoicing that by you the Lord feeds and clothes His poor, to whom
assuredly He could have given the possessions which He has bestowed on you, had He not in His
unspeakable mercy wished to justify them for their patient labour, and you for your works of
love.\par
{\smallskip\noindent\small ---Leo the Great, \textit{Sermo} XIX.3, tr. C.~L. Feltoe, NPNF, 2nd
ser., vol.~12 (New York, 1895).\par}
\end{quote}

\noindent He names the works of mercy that go with the fast: ``Let us defend the widow and serve the
orphan; let us comfort the afflicted and reconcile the estranged; let us take in the wanderer and
succour the oppressed; let us clothe the naked and cherish the sick'' (\textit{Sermo} 13, tr.
Bute, lesson vi). And he answers the fear of giving too much: ``Be stedfast, Christian giver: give
what you may receive, sow what you may reap, scatter what you may gather'' (\textit{Sermo} XVII.2,
tr. Feltoe). The simplest rule is Leo's own, in the sermon quoted in Section~\ref{sec:ember}:
spend in good deeds what the fast takes from indulgence. Give it before Christmas, and give some of it in person.
